% Metodologia -- fonte de dados, janela e seleção da amostra (parte do João).
% Vem antes das definições operacionais (Arthur) e da validação manual (Pedro).
% Números de dados/funil.csv, dados/repositorios.csv e do log da coleta.

\subsection{Fonte de dados e coleta}

Todos os dados vêm da API REST do GitHub, consultada por um \textit{pipeline} próprio em
Python que usa apenas a biblioteca \texttt{requests}. A coleta inteira é reproduzida por um
único comando. Cada resposta da API é gravada em um \textit{cache} local indexado pelo
\textit{hash} da URL, de modo que uma coleta interrompida continua de onde parou. A execução
final fez 11.989 chamadas à API e reaproveitou 58.459 respostas do \textit{cache}.

A busca de repositórios devolve no máximo 1.000 itens por consulta, e só a faixa de 1.000 a
1.500 estrelas tem 20.493 repositórios. Por isso, dividimos a busca em sete faixas de
estrelas (1.000--1.500, 1.500--2.500, 2.500--4.000, 4.000--7.000, 7.000--12.000,
12.000--25.000 e mais de 25.000). De cada faixa entraram os 1.000 repositórios mais
estrelados, num total de 7.000 candidatos. Os candidatos foram avaliados alternando as
faixas, para que a amostra não saísse inteira da faixa mais popular.

O cliente HTTP respeita a cota da API, esperando o \textit{reset} quando ela acaba, e repete
erros 5xx e falhas de rede com \textit{backoff} exponencial (1, 2, 4, 8 e 16~s), até seis
tentativas. Todas as listagens são paginadas pelo cabeçalho \texttt{Link}. A listagem de
\textit{workflow runs} também devolve no máximo 1.000 resultados, e o seu
\texttt{total\_count} satura em 2.500. Por isso, os \textit{runs} são pedidos mês a mês, e o
mês que passa do teto é dividido ao meio até caber nele.

\subsection{Janela de observação}

A janela tem 12 meses, de 01/10/2025 a 30/09/2026, inclusiva nas duas pontas, e é a mesma
para todos os repositórios. Só entram releases publicadas e \textit{workflow runs} criados
dentro dela, e só no \textit{default branch}. A idade do repositório é medida até o fim da
janela, e não até a data da execução, para que uma nova execução chegue ao mesmo valor.

\subsection{Seleção da amostra}

Os candidatos passaram por filtros em sequência. Cada repositório parou no primeiro filtro em
que falhou, e por isso tem um único motivo de descarte. A coleta parou ao completar 350
repositórios aprovados, acima do mínimo de 300. A Tabela~\ref{tab:funil} mostra o funil.

\begin{table}[ht]
\centering
\caption{Funil de seleção da amostra.}
\label{tab:funil}
\small
\begin{tabular}{lrrl}
\hline
\textbf{Etapa} & \textbf{Restantes} & \textbf{Descart.} & \textbf{Motivo do descarte} \\
\hline
Candidatos da busca (7 faixas) & 7.000 & -- & -- \\
Avaliados                      & 1.565 & 5.435 & não avaliados (amostra completa) \\
Usam GitHub Actions            & 1.057 & 508 & nenhum workflow \\
$\geq$ 5 releases na janela    & 439   & 618 & menos de 5 releases \\
\textit{Runs} mensuráveis      & 433   & 6   & $>$ 1.000 \textit{runs} em um dia \\
$\geq$ 50 \textit{runs} válidos & 352  & 81  & menos de 50 \textit{runs} válidos \\
Coleta sem erro                & 350   & 2   & erro definitivo da API \\
\hline
\textbf{Amostra final}         & \textbf{350} & -- & -- \\
\hline
\end{tabular}
\end{table}

Os 5.435 candidatos não avaliados nunca foram consultados, porque a coleta parou ao fechar
a amostra; eles não foram reprovados por critério algum. Os seis repositórios com mais de
1.000 \textit{runs} em um único dia foram descartados porque não há como recuperar todos os
\textit{runs} desses dias, nem saber quantos faltam. Mantê-los faria o CFR~(a) e o tempo de
recuperação serem calculados sobre dados truncados. Com isso, a amostra exclui os projetos de
maior volume de CI. Os dois erros definitivos foram um HTTP 504 recorrente e uma listagem de
releases que passou do limite de 10.000 resultados paginados da API.

A amostra final reúne 350 repositórios com 1.493 a 291.317 estrelas (mediana de 11.944),
concentrados nas faixas mais populares: 97 têm mais de 25.000 estrelas e 27 estão na faixa
de 1.000 a 1.500. As linguagens mais frequentes são TypeScript (73), Python (72), Go (45),
Rust (39) e JavaScript (28). A mediana é de 22 releases e 688 \textit{runs} válidos por
repositório na janela, e a idade mediana é de 6,0 anos.

A amostra não é aleatória. Dentro de cada faixa, só os 1.000 repositórios mais estrelados
chegaram a ser candidatos. Além disso, as faixas foram avaliadas na mesma proporção, mas os
repositórios mais populares passaram mais nos filtros, e por isso a amostra final pende para
eles. Os resultados valem para repositórios populares que usam GitHub Actions e publicam
releases, e não para o GitHub como um todo.

\subsection{Dicionário de dados}

O \textit{pipeline} gera dois arquivos. O \texttt{repositorios.csv} tem uma linha por
repositório avaliado, aprovado ou descartado, e o \texttt{funil.csv} tem uma linha por etapa
da seleção. O dicionário completo, com tipo, unidade e origem de cada coluna, está no
repositório do trabalho. As colunas de \texttt{repositorios.csv} formam cinco grupos:
(i) metadados (estrelas, linguagem, \textit{default branch}, idade e contribuidores);
(ii) contagens da seleção (releases na janela, sem rascunhos e pré-releases, e
\textit{runs} válidos, isto é, de \texttt{push} no \textit{default branch} que terminam em
sucesso ou falha); (iii) as métricas DORA em todas as variantes, com as contagens da
heurística de release corretiva e a proporção de episódios censurados; (iv) a nota de 1 a 4
por métrica e a categoria geral; e (v) o motivo do descarte. Valor vazio significa
\textit{não medido}, e não zero; métrica vazia não recebe nota e fica fora da mediana da
categoria geral.
